\chapter{Complex characters and induction}\label{chap:M07.21}
Throughout this chapter $G$ is a finite group, and representations and
vector spaces are over $\C$. Representations whose characters are taken
are finite-dimensional. We reuse the definitions, Schur lemmas and
Maschke theorem in Chapter~\ref{chap:M07.17}; their general-field versions
are not being silently strengthened to character classification.
The external D05 tensor, dual, trace and inner-product results cited by
their literal identifiers have their proofs in the D05 notes and the
declared foundation register of this project.

% Algebra, duals and averaging.
\KEntry{K07.21.group_algebra}{The group algebra and its modules}
The group algebra $\C[G]$ has basis $(e_g)_{g\in G}$, multiplication
$e_ge_h=e_{gh}$ extended bilinearly, and unit $e_1$.
We usually write $g$ for $e_g$, but the group identity then denotes
the algebra unit rather than the zero vector. For three basis elements,
$(e_ge_h)e_k=e_{(gh)k}=e_{g(hk)}=e_g(e_he_k)$; trilinearity
proves associativity on all elements. Scalars $ce_1$ commute with the
basis and hence with the whole algebra. Its dimension is $|G|$.

A representation $\rho$ extends uniquely to a unital algebra map
\[
 \C[G]\longrightarrow\End_\C(V),\qquad
 \sum_g a_g g\longmapsto\sum_g a_g\rho(g).
\]
Indeed multiplication of the two finite sums gives
$\sum_{g,h}a_gb_h\rho(gh)$, which is the required image of their
algebra product. Conversely a unital $\C[G]$-module gives a group
action, because the actions of $g$ and $g^{-1}$ are inverse. The
scalar action is the given complex vector-space action. These
constructions are inverse. A subspace is $G$-invariant exactly when it
is a submodule, and a linear map is equivariant exactly when it commutes
with the action of every algebra element: both assertions follow by
linear extension from the basis $G$. This proves the correspondence,
including the zero module and the case $G=1$, when $\C[G]=\C$.

\KEntry{K07.21.dual_tensor}{Dual and tensor representations}
For a representation on $V$, put
\[
 (g\phi)(v)=\phi(g^{-1}v)\qquad(\phi\in V^*).
\]
This is linear in $\phi$, and
$g(h\phi)(v)=\phi(h^{-1}g^{-1}v)=((gh)\phi)(v)$.
The identity acts identically and the action of $g^{-1}$ is inverse.
Thus it defines the dual representation; omitting the inverse would
reverse the product law. In dual bases its matrix is
$(\rho(g^{-1}))^{\mathsf T}$, so
\[
 \chi_{V^*}(g)=\chi_V(g^{-1}).
\]
Finite-dimensional duality also preserves irreducibility. If
$0<U<V^*$ were invariant, its annihilator in $V$ would be a proper
nonzero invariant subspace: for $v$ annihilating $U$ and $\phi\in U$,
$\phi(gv)=(g^{-1}\phi)(v)=0$. Its positive, proper dimension follows
from the D05 annihilator dimension formula
\texttt{K05.07.annihilator\_dimension}. This contradicts irreducibility
of $V$. Conversely apply the same assertion to $V^*$ and the natural
equivariant identification $V\simeq V^{**}$, whose equivariance follows
by evaluating at every $\phi$. Zero spaces remain zero and are not
irreducible.

For representations $V,W$, define
$g(v\otimes w)=gv\otimes gw$. The tensor universal property and
functoriality, \texttt{K05.15.tensor\_universal} and
\texttt{K05.15.tensor\_maps}, give linear operators satisfying the
group law on elementary tensors and therefore everywhere. In product
bases the diagonal entries are $\rho(g)_{ii}\sigma(g)_{jj}$; summing
gives
\[
 \chi_{V\otimes W}(g)=\chi_V(g)\chi_W(g).
\]
If one space is zero both sides are zero. These are representation
structures on the D05 tensor product, not a new tensor construction.

\KEntry{K07.21.square_characters}{Symmetric and exterior square characters}
The flip $F(v\otimes w)=w\otimes v$ on $V\otimes V$ commutes
with $A\otimes A$ for every linear $A$. The two projections
$P_+=(I+F)/2$ and $P_-=(I-F)/2$ have complementary images,
identified with $\operatorname{Sym}^2V$ and $\Lambda^2V$ by
\texttt{K05.15.symmetrizers}, \texttt{K05.15.symmetric} and
\texttt{K05.15.exterior}. Consequently both are invariant for the
tensor representation. If a map $T$ commutes with a projection $P$,
then $T$ is block diagonal on $\im P\oplus\ker P$, while $TP$
has blocks $T|_{\im P},0$; thus $\Tr(TP)=\Tr(T|_{\im P})$.
For a matrix $A$, product-basis entries give
\[
 \Tr((A\otimes A)F)=\sum_{i,j}A_{ij}A_{ji}=\Tr(A^2).
\]
Taking $A=\rho(g)$ therefore proves
\[
 \chi_{\operatorname{Sym}^2 V}(g)
  =\frac{\chi_V(g)^2+\chi_V(g^2)}2,\qquad
 \chi_{\Lambda^2 V}(g)
  =\frac{\chi_V(g)^2-\chi_V(g^2)}2.
\]
The argument includes dimensions zero and one. Division by two is valid
over $\C$; these projection formulas are not asserted in characteristic
two.

\KEntry{K07.21.unitary_averaging}{Invariant Hermitian metrics}
Choose any positive definite Hermitian form $h_0$ on a finite-dimensional
$V$, using a basis. Fix linearity in the first argument and conjugate
linearity in the second; \texttt{K05.12.convention\_bridge} relates this
to the opposite convention. Set
\[
 h(v,w)=\frac1{|G|}\sum_{g\in G}h_0(gv,gw).
\]
The form is Hermitian and sesquilinear by addition. If $v\ne0$, every
$gv$ is nonzero, so the summands $h_0(gv,gv)$ are positive and
$h(v,v)>0$. Reindexing $g\mapsto gk$ proves
$h(kv,kw)=h(v,w)$. Gram--Schmidt from
\texttt{K05.12.gram\_schmidt} gives an orthonormal basis in which every
$\rho(g)$ is unitary. The zero space has its unique Hermitian form,
which is positive definite with the usual vacuous nonzero-vector test.
Since a unitary matrix has inverse its conjugate transpose,
\[
 \chi_V(g^{-1})=\overline{\chi_V(g)},\qquad
 \chi_{V^*}(g)=\overline{\chi_V(g)}.
\]
The first identity concerns finite complex representations of a finite
group; it does not follow from trace invariance alone over an arbitrary
field.

\KEntry{K07.21.hom_pairing}{Averaging Hom spaces and character inner products}
For functions on $G$, define
\[
 \langle f,k\rangle_G=\frac1{|G|}\sum_{g\in G}
 f(g)\overline{k(g)}.
\]
This is a positive definite Hermitian inner product: the squared norm is
the average of $|f(g)|^2$, and vanishes exactly when every value vanishes.
It is linear in its first argument. On $\Hom_\C(V,W)$ use the action
\[
 gT=\sigma(g)T\rho(g)^{-1}.
\]
Composition of these operators proves the group law, and its invariant
subspace is exactly $\Hom_G(V,W)$.
The average $Q=|G|^{-1}\sum_g(g\cdot-)$ has invariant image by
reindexing and is the identity on invariants. Hence $Q^2=Q$ and
$\im Q=\Hom_G(V,W)$. Its decomposition into image and kernel
shows $\Tr Q=\dim_\C\Hom_G(V,W)$, since its matrix there is
$\operatorname{diag}(I,0)$.

In bases of $V,W$, the trace of the map $T\mapsto BT A$ on the
matrix space is $\Tr(B)\Tr(A)$: the coefficient of a matrix unit
$E_{ba}$ in its own image is $B_{bb}A_{aa}$, and these units form
a basis. Take $B=\sigma(g),A=\rho(g)^{-1}$ and use the preceding
unitary character identity. Trace linearity gives the exact formula
\begin{equation}\label{eq:character-hom-pairing}
 \langle\chi_W,\chi_V\rangle_G
   =\dim_\C\Hom_G(V,W).
\end{equation}
The integer on the right is viewed in $\C$. This also proves that these
particular inner products are real nonnegative integers. When either
space is zero both sides vanish.

\KEntry{K07.21.matrix_orthogonality}{Schur orthogonality of matrix coefficients}
Choose pairwise inequivalent irreducible representations $V_i$, with
unitary bases, and write $d_i=\dim V_i>0$. No completeness of the
chosen family is assumed yet. Then
\begin{equation}\label{eq:matrix-orthogonality}
 \frac1{|G|}\sum_{g\in G}\rho_i(g)_{ab}
      \overline{\rho_j(g)_{cd}}
 =\begin{cases}\delta_{ac}\delta_{bd}/d_i&i=j,\\0&i\ne j.
 \end{cases}
\end{equation}
To prove it, average a map $T:V_j\to V_i$ as in the preceding entry.
The result $Q(T)$ is an intertwiner. If $i\ne j$, Schur's
zero-or-isomorphism lemma makes it zero. If $i=j$, scalar Schur gives
$Q(T)=\lambda I$. Trace invariance under conjugation gives
$d_i\lambda=\Tr Q(T)=\Tr T$, so $Q(T)=(\Tr T/d_i)I$.
Now take $T$ with its single nonzero entry $T_{bd}=1$. The $(a,c)$
entry of the average is
$|G|^{-1}\sum_g\rho_i(g)_{ab}\rho_j(g)^{-1}_{dc}$.
Unitarity identifies the last entry with
$\overline{\rho_j(g)_{cd}}$, and the scalar answer has entry
$\delta_{bd}\delta_{ac}/d_i$, as required.

Summing the diagonal coefficient pairings proves
\begin{equation}\label{eq:character-row-orthogonality}
 \langle\chi_i,\chi_j\rangle_G=\delta_{ij}.
\end{equation}
This also follows from \eqref{eq:character-hom-pairing} and the two
Schur lemmas. All matrix coefficients in this family are nonzero
orthogonal functions after the indicated normalization, not only their
diagonal sums.

% Regular representations and completeness.
\KEntry{K07.21.regular_decomposition}{All irreducibles and the regular representation}
\begin{theorem}\label{thm:regular-complex-decomposition}
There are finitely many irreducible complex representations of $G$ up
to equivalence. Choose one of each, $V_1,\ldots,V_r$, and put
$d_i=\dim V_i$. Then
\[
 \C[G]\cong_G\bigoplus_{i=1}^r V_i^{\oplus d_i},
 \qquad |G|=\sum_{i=1}^r d_i^2.
\]
Every irreducible representation of a finite group is finite-dimensional,
even if this was not assumed in its definition.
\end{theorem}
\begin{proof}
For an irreducible $V$ choose $v\ne0$. The span of the orbit $Gv$
is nonzero and invariant, hence is $V$. The map
$\C[G]\to V$, $\sum a_gg\mapsto\sum a_g gv$, is therefore
surjective and equivariant. It also proves $\dim V\leq |G|$.
By Maschke and \ref{K07.17.complete_reducibility}, the finite-dimensional
regular representation is a finite direct sum of irreducibles.
The displayed surjection is nonzero on at least one summand; its
restriction there is an equivalence by Schur. Thus its finitely many
summand types are all irreducible types, proving finiteness.

Collect equal types to write the regular module as
$\bigoplus_i V_i^{\oplus m_i}$. Left multiplication by $g$ permutes
the basis $G$ and has a fixed basis vector exactly when $g=1$.
Thus its character is $|G|$ at $1$ and zero elsewhere. Pairing that
character with $\chi_i$ and applying row orthogonality gives
\[
 m_i=\langle\chi_{\C[G]},\chi_i\rangle_G
 =\overline{\chi_i(1)}=d_i.
\]
Taking dimensions proves the square sum. In particular every chosen
type really occurs. There is at least one type, since the nonzero
regular module contains the trivial line spanned by $\sum_g g$.
For $G=1$, the result is the single one-dimensional trivial type.
\end{proof}
There are $\sum_i d_i^2=|G|$ matrix coefficient functions in
\eqref{eq:matrix-orthogonality}. They are independent and hence form
an orthogonal basis of the entire $|G|$-dimensional function space
$\C^G$. This completeness is now justified by the regular module,
not presumed while proving orthogonality.

\KEntry{K07.21.character_completeness}{The class count and complete character basis}
\begin{theorem}\label{thm:character-completeness}
The number $r$ of irreducible complex representation types equals the
number of conjugacy classes of $G$. Their characters form an orthonormal
basis of the space of complex class functions on $G$.
\end{theorem}
\begin{proof}
The representation actions extend to an algebra map
\[
 \Phi:\C[G]\longrightarrow\prod_{i=1}^r\End_\C(V_i).
\]
If $a$ is in its kernel, it acts by zero on every irreducible and
therefore on every summand of the regular module. Left multiplication
by $a$ on $\C[G]$ is zero, so evaluating it at the algebra unit gives
$a=0$. Thus $\Phi$ is injective. By the square sum, both vector
spaces have dimension $|G|$, so it is surjective and hence an algebra
isomorphism. This proves the needed special group-algebra decomposition
without assuming a general structure theorem for associative algebras.

An endomorphism commuting with all matrix units is scalar: commuting
with each diagonal matrix unit makes its off-diagonal entries zero;
commuting with $E_{ab}$ then equates every two diagonal entries.
This includes dimension one. Hence the center of the product algebra
has dimension $r$, one scalar per factor, and so does $Z(\C[G])$.
On the other hand, $a=\sum_g a_g g$ is central exactly when
$hah^{-1}=a$ for every $h\in G$, equivalently when its coefficients
are constant on conjugacy classes. The sums of the basis elements in
each class form a basis of this center: their supports are disjoint,
and every constant-on-class coefficient list is their linear combination.
Thus its dimension is the number $k$ of classes, giving $r=k$.

Class functions form a $k$-dimensional vector space with the class
indicator basis. The $r=k$ irreducible characters lie in this space by
\ref{K07.17.characters} and are orthonormal by
\eqref{eq:character-row-orthogonality}. They are consequently a basis.
\end{proof}

\KEntry{K07.21.multiplicity_classification}{Multiplicities and equality of characters}
For a finite-dimensional representation $V$, complete reducibility
gives $V\cong\bigoplus_i V_i^{\oplus m_i}$ with $m_i\geq0$.
Taking characters and pairing with $\chi_j$ proves
\[
 m_j=\langle\chi_V,\chi_j\rangle_G,
 \qquad \langle\chi_V,\chi_V\rangle_G=\sum_i m_i^2.
\]
Thus the multiplicities are unique. Two finite-dimensional complex
representations are equivalent if and only if their characters are equal:
equal characters give the same multiplicities and isomorphisms with the
same direct sum, while an equivalence preserves trace. A representation
is irreducible exactly when its character has squared norm one, because
a sum of squares of nonnegative integers is one exactly when one is
one and the others are zero. The zero space has all multiplicities and
character zero and has squared norm zero; it is not an exception to the
equivalence criterion. More generally
\[
 \dim\Hom_G(V,W)=\sum_i m_i(V)m_i(W)
\]
follows either by the character inner product or by Schur on summands.
These conclusions have the complex, finite-group and finite-dimensional
hypotheses fixed above; the modular counterexample in
\ref{K07.17.characters} prevents dropping them indiscriminately.

\KEntry{K07.21.column_orthogonality}{The second orthogonality relation}
Choose class representatives $g_1,\ldots,g_r$, and write $c_j$ for
the size of the class of $g_j$. The square matrix
$U_{ij}=\sqrt{c_j/|G|}\,\chi_i(g_j)$ has orthonormal rows by the
first orthogonality relation. Hence $UU^*=I$; finite-dimensional
invertibility then gives $U^*U=I$. Comparing its $(j,l)$ entries and
conjugating if necessary proves
\[
 \sum_i\chi_i(g)\overline{\chi_i(h)}
 =\begin{cases}|C_G(g)|&g,h\text{ are conjugate},\\0&\text{otherwise}.
 \end{cases}
\]
Indeed at equal class indices the column norm is $|G|/c_j$, which is
$|C_G(g_j)|$ by conjugacy orbit--stabilizer
\ref{K07.05.conjugacy_counting}; unequal indices give zero.
Characters are constant on classes, so this covers arbitrary $g,h$.
The square-matrix step uses the class-count theorem: row orthogonality
of a rectangular partial character table would not imply this conclusion.

\KEntry{K07.21.s3_table}{The complete character table of the symmetric group on three letters}
The three classes in $S_3$ have representatives $1,s=(12),a=(123)$
and sizes $1,3,2$, by cycle-type conjugacy
\ref{K07.02.conjugacy_cycle_type}. The trivial and sign characters
are $(1,1,1)$ and $(1,-1,1)$. On $\C^3$, the permutation action
preserves $W=\{(x_1,x_2,x_3):x_1+x_2+x_3=0\}$ and the trivial
line $L=\C(1,1,1)$. Every vector is uniquely the sum of its mean
coordinate vector in $L$ and a vector in $W$, so
$\C^3=L\oplus W$. The permutation character counts fixed basis
points and has values $(3,1,0)$. Therefore $\chi_W=(2,0,-1)$.
Its squared norm is $(4+3\cdot0+2\cdot1)/6=1$, proving
irreducibility by the multiplicity criterion. It differs from the two
one-dimensional characters. Their squared dimensions sum to
$1+1+4=6$, so the regular square sum excludes another irreducible type.
The complete table is thus
\[
\begin{array}{c|rrr}
 &1&(12)&(123)\\ \hline
\text{class size}&1&3&2\\
\chi_{\mathrm{triv}}&1&1&1\\
\chi_{\mathrm{sgn}}&1&-1&1\\
\chi_W&2&0&-1
\end{array}
\]
The representation itself is explicit. In the basis $e_1-e_3,e_2-e_3$
of $W$, the matrices of $s$ and $a$ are respectively
\[
 \begin{pmatrix}0&1\\1&0\end{pmatrix},\qquad
 \begin{pmatrix}-1&-1\\1&0\end{pmatrix}.
\]
Their columns are the coordinates of the permuted basis vectors, so
their trace values and the common representation law require no
guess from a numerical character table.

% Restriction, induction and reciprocity.
\KEntry{K07.21.induction_coset_model}{The induced representation and its coset model}
Restriction $\operatorname{Res}_H^G W$ for $H\leq G$ is the same
vector space with the action restricted to $H$; its group laws are
inherited. To construct induction from an $H$-representation $V$, form
the ordinary complex tensor product $\C[G]\otimes_\C V$ from D05,
and quotient by the subspace generated by
\[
 gh\otimes v-g\otimes hv\qquad(g\in G,h\in H,v\in V).
\]
Denote a quotient generator by $[g,v]$, and denote the quotient by
$\operatorname{Ind}_H^G V$. This explicit construction is also written
$\C[G]\otimes_{\C[H]}V$; no unproved module-tensor construction is
being presumed by that notation. Left multiplication on the first
factor preserves the relation subspace, since $k$ sends the displayed
relation to $kgh\otimes v-kg\otimes hv$. It therefore defines
$k[g,v]=[kg,v]$ on the quotient, with inverse given by $k^{-1}$.
This is a $G$-representation.

Choose representatives $T$ for the left cosets $G/H$. There is a
vector-space isomorphism
\begin{equation}\label{eq:induction-coset-sum}
 \bigoplus_{t\in T}V\longrightarrow\operatorname{Ind}_H^G V,
 \qquad(v_t)_t\longmapsto\sum_t[t,v_t].
\end{equation}
For its inverse, write each $g$ uniquely as $th$ and send $g\otimes v$
to the vector with $hv$ in its $t$-component and zero elsewhere.
Extend linearly over the basis $G$ and linearly in $v$. It kills the
relations: for $h_0\in H$, the two expressions $gh_0\otimes v$
and $g\otimes h_0v$ both give $hh_0v$ in that component.
Thus it descends to the quotient. Its composite with
\eqref{eq:induction-coset-sum} is the identity. In the other order,
$[g,v]=[th,v]=[t,hv]$, proving identity on spanning generators.

If $kt=t'h$, the action in this model sends the $t$-copy of $v$ to
the $t'$-copy of $hv$. This is the explicit coset block-matrix model.
For finite-dimensional $V$ it proves
$\dim\operatorname{Ind}_H^G V=[G:H]\dim V$. Different representative
choices give equivalent coordinate models of the same constructed
quotient. Induction sends the zero space to zero. For $H=G$ the map
$[g,v]\mapsto gv$ is an equivalence with inverse $v\mapsto[1,v]$;
for $H=1$, induction is the regular representation tensored with $V$.
If $V\ne0$ is finite-dimensional and $H<G$, its larger dimension shows that induction
does not mean an extension of the original action on the same space.

\KEntry{K07.21.induced_character}{The induced-character formula}
\begin{theorem}\label{thm:induced-character-formula}
For $H\leq G$, an $H$-character $\psi$, and left-coset representatives
$T$, the character of the induced representation is
\[
 \operatorname{Ind}_H^G\psi(g)
 =\sum_{\substack{t\in T\\t^{-1}gt\in H}}\psi(t^{-1}gt)
 =\frac1{|H|}\sum_{\substack{x\in G\\x^{-1}gx\in H}}
       \psi(x^{-1}gx).
\]
In particular it vanishes if $g$ is not conjugate to an element of $H$,
and its value at $1$ is $[G:H]\psi(1)$.
\end{theorem}
\begin{proof}
In the coset model, the action of $g$ sends the $t$-block to a different
block unless $gtH=tH$. A block moved elsewhere contributes zero
to the full matrix diagonal. On a fixed block its matrix is the action
of $t^{-1}gt\in H$ on $V$, so its trace is the first summand.
This proves the first formula. For $x=th$, the condition on $x$ is
equivalent to the condition on $t$; when it holds,
$x^{-1}gx=h^{-1}(t^{-1}gt)h$. The character is constant on
$H$-conjugacy classes, so all $|H|$ choices of $h$ give the same
value. Sum coset by coset to get the second formula. The two final
assertions follow by an empty sum or by all cosets being fixed,
respectively. For $V=0$ the formula is identically zero.
\end{proof}
The second expression defines induction of an arbitrary complex class
function as well: replacing $g$ by $kgk^{-1}$ and $x$ by $kx$
leaves the sum unchanged. This observation will be used only where its
class-function hypothesis is stated.

\KEntry{K07.21.frobenius_reciprocity}{Frobenius reciprocity, as maps and as characters}
\begin{theorem}\label{thm:frobenius-reciprocity}
There is a natural linear isomorphism
\[
 \Hom_G(\operatorname{Ind}_H^G V,W)
 \simeq\Hom_H(V,\operatorname{Res}_H^G W).
\]
For class functions $f$ on $H$ and $k$ on $G$,
\[
 \langle\operatorname{Ind}_H^G f,k\rangle_G
   =\langle f,\operatorname{Res}_H^G k\rangle_H.
\]
In particular these formulas apply to characters and give the
multiplicity of an irreducible $G$-representation in an induced one.
\end{theorem}
\begin{proof}
Send $T$ to $A(v)=T[1,v]$. For $h\in H$,
$A(hv)=T[1,hv]=T[h,v]=hT[1,v]$, proving $H$-equivariance.
Conversely, for an $H$-equivariant $A$, define
$T[g,v]=gA(v)$. The basis $G$ and the tensor universal property
extend this formula linearly before quotienting. It kills the relation
because $ghA(v)=gA(hv)$, so it descends. It is $G$-equivariant since
$T[k g,v]=kgA(v)=kT[g,v]$. These constructions are inverse:
evaluation at $[1,v]$ recovers $A$, while $G$-equivariance forces
$T[g,v]=gT[1,v]$ on all spanning generators. They are linear,
including when either space is zero.

For the class-function formula substitute the induced sum into the
inner product. All sums are finite. Write $g=xhx^{-1}$ in its inner
sum; for a fixed $x$ this is a bijection between eligible $g$ and
$h\in H$. Since $k$ is a class function, the result is
\[
 \begin{aligned}
 \frac1{|G||H|}\sum_{x\in G}\sum_{h\in H}
       f(h)\overline{k(xhx^{-1})}
 &=\frac1{|G||H|}\sum_{x\in G}\sum_{h\in H}
       f(h)\overline{k(h)}\\
 &=\frac1{|H|}\sum_{h\in H}f(h)\overline{k(h)}.
 \end{aligned}
\]
This is the required $H$-inner product. For characters the multiplicity
interpretation follows from the earlier classification formula, so no
assertion that an arbitrary induced representation is irreducible is
needed.
\end{proof}

\KEntry{K07.21.induction_transitivity}{Transitivity and additivity of induction}
For $H\leq K\leq G$ there is an explicit equivalence
\[
 \operatorname{Ind}_K^G(\operatorname{Ind}_H^K V)
 \simeq_G\operatorname{Ind}_H^G V,
 \qquad [g,[k,v]]\longmapsto[gk,v].
\]
To verify descent, the inner relation $[kh,v]=[k,hv]$ is taken to
$[gkh,v]=[gk,hv]$, and the outer $K$-relation
$[gk_0,[k,v]]=[g,k_0[k,v]]$ is taken on both sides to
$[gk_0k,v]$. Tensor linearity and quotient universality from D05
therefore construct the map. Its inverse sends $[g,v]$ to
$[g,[1,v]]$: the $H$-relation is respected because
\[
 [gh,[1,v]]=[g,h[1,v]]=[g,[h,v]]=[g,[1,hv]].
\]
One composite fixes $[g,v]$. The other sends $[g,[k,v]]$ to
$[gk,[1,v]]=[g,[k,v]]$. Spanning proves both are identities,
and both commute with the left $G$-action. This also treats equal
adjacent subgroups and zero spaces.

Induction preserves direct sums: the map
$[g,(v,w)]\mapsto([g,v],[g,w])$ descends because each component
respects the same $H$-relations. Its inverse is the sum of the induced
inclusions, and their composites are identities on generators. Thus
\[
 \operatorname{Ind}_H^G(\psi+\theta)
 =\operatorname{Ind}_H^G\psi+\operatorname{Ind}_H^G\theta,
 \qquad
 \operatorname{Ind}_K^G\operatorname{Ind}_H^K\psi
 =\operatorname{Ind}_H^G\psi.
\]
These are proved representation equivalences before being read as
character identities; no unproved associativity of module tensors is
required.

\KEntry{K07.21.induction_examples}{Induction in the symmetric group}
For any $H\leq G$, induction of its one-dimensional trivial
representation is the permutation representation on $G/H$: the coset
model has one basis vector per coset and each $h$ acts on its coefficient
as the identity. In particular induction from the trivial subgroup gives
the regular representation. If $|G|>1$, this is reducible: its invariant
line spanned by the sum of all basis vectors is nonzero and proper.
Thus inducing an irreducible representation need not preserve
irreducibility.

For $G=S_3$ and $H=\langle(12)\rangle$, the coset set is identified
with $\{1,2,3\}$ by $gH\mapsto g(3)$: $H$ is exactly the
stabilizer of $3$, so the orbit--coset theorem proves bijectivity and
equivariance. Hence induction of the trivial character is
$\chi_{\mathrm{triv}}+\chi_W$. If $\varepsilon$ is the nontrivial
character of $H$, the induced character has values $(3,-1,0)$.
Its degree is three; a transposition fixes exactly one coset and acts
on that block by the nonidentity element of $H$, of value $-1$;
a three-cycle fixes no coset. Therefore it equals
$\chi_{\mathrm{sgn}}+\chi_W$, and character classification proves
the corresponding representation equivalence.

Finally take $H=A_3=\langle a\rangle$, choose a primitive cube root
$\omega$, and let $\psi(a)=\omega$. The representatives $1,s$,
with $s=(12)$, give induced values $2$ at $1$, zero at a
transposition, and $\omega+\omega^2=-1$ at $a$, since
$s^{-1}as=a^{-1}$. Thus $\operatorname{Ind}_{A_3}^{S_3}\psi$
has character $\chi_W$ and is irreducible. The distinct conjugate
character $\overline\psi$ induces the same representation, showing
induction is not injective on representation types. Nor does $\psi$
extend to a one-dimensional representation of $S_3$: such a map has
abelian target and kills the commutator subgroup $A_3$, computed in
\ref{K07.18.nilpotence_examples}, whereas $\psi(a)\ne1$.
This distinguishes an induced representation from an extension and
concludes the bounded induction theory without an assertion about
general Frobenius groups or classification of simple groups.
